\section{Preliminaries and Problem Setup}

We work in the Lorentz model because it gives closed-form geodesics and parallel transport. For $K>0$, the $d$-dimensional hyperbolic space of sectional curvature $-K$ is
\begin{equation}
\begin{gathered}
\HH_{-K}^{d}=\{x\in\RR^{d+1}:\inner{x}{x}=-1/K,\;x_0>0\},\\
\inner{x}{y}=-x_0y_0+\sum_{i=1}^{d}x_iy_i .
\end{gathered}
\end{equation}
The tangent space at $x$ is $T_x\HH_{-K}^{d}=\{v:\inner{x}{v}=0\}$ with norm $\|v\|_x=\sqrt{\inner{v}{v}}$. The exponential map is
\begin{equation}
\Exp_x^K(v)=\cosh(\sqrt{K}r)x+\frac{\sinh(\sqrt{K}r)}{\sqrt{K}r}v,\qquad r=\|v\|_x,
\label{eq:exp-main}
\end{equation}
with the usual convention at $r=0$. This formula shows why the product $\sqrt{K}r$ is central: it is the argument of the hyperbolic sine and cosine terms that control both useful volume growth and harmful numerical amplification.

For momentum, the previous tangent vector must be moved before it can be combined with the current gradient. Along the geodesic from $x$ to $y$ at fixed curvature, the Lorentz parallel transport of $v\in T_x\HH^d_{-K}$ is
\begin{equation}
\PT_{x\to y}^{K}(v)=v+\frac{K\inner{y}{v}}{1-K\inner{x}{y}}(x+y),
\label{eq:pt-main}
\end{equation}
which lies in $T_y\HH^d_{-K}$ and preserves the Lorentz norm. When curvature changes from $K$ to $K'$, the natural Lorentz rescaling maps points and tangent vectors as $\Psi_{K\to K'}(x)=\sqrt{K/K'}x$ and $\mathcal{C}_{K\to K'}(v)=\sqrt{K/K'}v$. These maps preserve the hyperboloid constraint and give the first-order cross-curvature consistency used by the optimizer.

The learning problem has $G$ structural groups. Each group $g$ has a parameter point $x_{g,t}\in\HH^{d_g}_{-K_{g,t}}$, a momentum vector $m_{g,t}\in T_{x_{g,t}}\HH^{d_g}_{-K_{g,t}}$, and a log-curvature coordinate $u_{g,t}=\log K_{g,t}$ implemented with a positive softplus parameterization and clipping to $[K_{\min},K_{\max}]$. The joint objective is
\begin{equation}
\mathcal{J}(x,u)=F(x,e^u)+\lambda_u\sum_g (u_g-u_{\mathrm{ref}})^2.
\label{eq:joint-main}
\end{equation}
The regularizer is not a claim that all curvatures should equal the reference; it is a numerical guard against runaway curvature. In the KG experiments, WN18RR reaches a boundary-like regime near $K=2$ for the highest-curvature relation, whereas synthetic runs often remain in an interior regime. The distinction matters for theory and is stated explicitly in Section~\ref{sec:theory}.

The experimental method labels are fixed throughout the paper. B3 is fixed-curvature Riemannian momentum with $K=1$. B4 adds grouped learnable curvature and maps points as curvature changes. B5 is the full HyperbolicMomentum variant with learnable curvature, curvature-normalized geodesic steps, and cross-curvature momentum transport. In shared-entity KG architectures, relation-specific curvature is used at scoring time while entity embeddings are shared across relations; therefore B5's cross-curvature momentum mechanism has no coherent home manifold and the KG comparison is interpreted as B4 versus B3, not as evidence for B5 beating B4.
